\documentclass[11pt]{article}

\usepackage[T1]{fontenc}
\usepackage[utf8]{inputenc}
\usepackage{lmodern}
\usepackage[a4paper,margin=2.6cm]{geometry}
\usepackage{amsmath,amssymb,amsthm,mathtools}
\usepackage{booktabs}
\usepackage{microtype}
\usepackage{hyperref}

\hypersetup{hidelinks}
\allowdisplaybreaks

\newtheorem{proposition}{Proposition}[section]
\theoremstyle{definition}
\newtheorem{definition}[proposition]{Definition}
\newtheorem{example}[proposition]{Example}
\theoremstyle{remark}
\newtheorem{remark}[proposition]{Remark}

\newcommand{\R}{\mathbb{R}}
\newcommand{\E}{\mathbb{E}}
\newcommand{\eps}{\epsilon}
\newcommand{\one}{\mathbf{1}}
\newcommand{\al}{\alpha}
\newcommand{\be}{\beta}
\newcommand{\Xx}{\mathcal{X}}
\newcommand{\Yy}{\mathcal{Y}}
\newcommand{\Couplings}{\mathcal{U}}
\newcommand{\MK}{\mathcal{L}}
\newcommand{\Gaussian}{\mathcal{N}}
\newcommand{\Dd}{\mathcal{D}}
\newcommand{\Ee}{\mathcal{E}}
\newcommand{\Ss}{\mathcal{S}}
\newcommand{\dd}{\,\mathrm{d}}
\newcommand{\ip}[2]{\left\langle #1,#2\right\rangle}
\newcommand{\norm}[1]{\left\lVert #1\right\rVert}
\newcommand{\opnorm}[1]{\left\lVert #1\right\rVert_{\mathrm{op}}}

\DeclareMathOperator{\KL}{KL}
\DeclareMathOperator{\Cov}{Cov}
\DeclareMathOperator{\Var}{Var}
\DeclareMathOperator{\essinf}{ess\,inf}
\DeclareMathOperator{\esssup}{ess\,sup}

\title{Hessian Geometry, Local Rates, and Acceleration of Sinkhorn}
\author{}
\date{23 July 2026}

\begin{document}

\maketitle

\begin{abstract}
This note isolates one conditional-expectation operator that controls
three aspects of entropic optimal transport: dual curvature, the local
convergence of Sinkhorn's algorithm, and semi-dual conditioning. At the
optimal coupling $\pi_\eps^\star$, write
$T_\eps:=T_{\pi_\eps^\star}$. The positive full-dual curvature operator is
$\eps^{-1}\left(\begin{smallmatrix}I&T_\eps\\T_\eps^\ast&I\end{smallmatrix}\right)$,
the Jacobian of one complete Sinkhorn cycle is $T_\eps^\ast T_\eps$, and exact
elimination of one potential gives the semi-dual Hessian
$\eps^{-1}(I-T_\eps^\ast T_\eps)$. This common structure yields the local rate,
strict-gap criteria, small-$\eps$ asymptotics, and explicit parameters
for blockwise over-relaxation. It also explains why variable
projection and quasi-Newton methods are natural alternatives to long
scaling iterations.
\end{abstract}

\section{Dual Hessian and conditional expectation}

The correct geometry for the dual potentials is not an unweighted
Euclidean space: it is the product of the two marginal $L^2$ spaces. In
this geometry, the Hessian has an especially simple block form.

\paragraph{Entropic primal and dual.}
Let $\Xx$ and $\Yy$ be standard Borel spaces, let
$\al\in\mathcal P(\Xx)$ and $\be\in\mathcal P(\Yy)$, and let
$c:\Xx\times\Yy\to\R\cup\{+\infty\}$ be measurable. The
KL-regularized transport value is
\begin{equation}
  \label{eq:primal}
  \MK_c^\eps(\al,\be)
  :=
  \inf_{\pi\in\Couplings(\al,\be)}
  \left\{
    \int c\,\dd\pi
    +
    \eps\KL(\pi\,\|\,\al\otimes\be)
  \right\}.
\end{equation}
This is the entropy-regularized transport problem introduced
computationally by Cuturi \cite{Cuturi2013}. Its concave dual objective
is
\begin{equation}
  \label{eq:dual-energy}
  \Dd_\eps(f,g)
  :=
  \int_{\Xx}f\,\dd\al+\int_{\Yy}g\,\dd\be
  -\eps\int_{\Xx\times\Yy}
  \left[
  \exp\left(\frac{f(x)+g(y)-c(x,y)}{\eps}\right)-1
  \right]
  \dd\al(x)\dd\be(y)
\end{equation}
and $\MK_c^\eps(\al,\be)=\sup_{f,g}\Dd_\eps(f,g)$ under the standard
duality assumptions. For any potentials for which the integral is
finite, denote their Gibbs measure by
\begin{equation}
  \label{eq:gibbs-measure}
  \dd\pi_{f,g}(x,y)
  =
  \exp\left(\frac{f(x)+g(y)-c(x,y)}{\eps}\right)
  \dd\al(x)\dd\be(y)
\end{equation}
and write $(\pi_{f,g})_1,(\pi_{f,g})_2$ for its marginals. The first
variation is
\begin{equation}
  \label{eq:dual-first-variation}
  D\Dd_\eps(f,g)[h,k]
  =
  \int_{\Xx}h\,\dd\bigl(\al-(\pi_{f,g})_1\bigr)
  +
  \int_{\Yy}k\,\dd\bigl(\be-(\pi_{f,g})_2\bigr).
\end{equation}
Consequently, $\pi_{f,g}\in\Couplings(\al,\be)$ exactly when $(f,g)$
is stationary. Since $\Dd_\eps$ is concave, every such pair is a dual
maximizer. Assume henceforth that a maximizer $(f^\star,g^\star)$
exists, and set
\begin{equation}
  \label{eq:optimal-gibbs-plan}
  \pi_\eps^\star:=\pi_{f^\star,g^\star}
  \in\Couplings(\al,\be).
\end{equation}
The gauge transformation
$(f,g)\mapsto(f+s,g-s)$ leaves \eqref{eq:dual-energy} unchanged. We
therefore equip the potentials with the quotient Hilbert space
\begin{equation}
  \label{eq:dual-quotient-space}
  \mathcal H_{\mathrm q}
  :=
  \frac{L^2(\al)\oplus L^2(\be)}
       {\operatorname{span}\{(\one,-\one)\}},
  \qquad
  \norm{[(h,k)]}_{\mathcal H_{\mathrm q}}^2
  :=
  \inf_{s\in\R}
  \left\{\norm{h+s\one}_\al^2+\norm{k-s\one}_\be^2\right\}.
\end{equation}
Each class has a unique representative orthogonal to the gauge
direction, characterized by
$\int h\dd\al=\int k\dd\be$. Notice that this condition does not center
both potentials: the common-constant direction $(\one,\one)$ remains
in the quotient.

We also use the centered spaces
\[
  L_0^2(\al):=\left\{h:\int h\,\dd\al=0\right\},
  \qquad
  L_0^2(\be):=\left\{k:\int k\,\dd\be=0\right\},
\]
endowed with their usual weighted norms. We write
\begin{equation}
  \label{eq:centering-projectors}
  \operatorname{cent}_{\al}h:=h-\int h\dd\al,
  \qquad
  \operatorname{cent}_{\be}k:=k-\int k\dd\be
\end{equation}
for the orthogonal centering projectors.
The product $L_0^2(\al)\oplus L_0^2(\be)$ isolates the nonconstant
singular modes, but it has codimension two and is not a representative
of the quotient \eqref{eq:dual-quotient-space}.

\paragraph{Coupling operator.}
The operator governing the local theory is meaningful for every
coupling, independently of entropic optimality.

\begin{definition}[Conditional operator and maximal correlation]
  \label{def:coupling-operator}
  For $\pi\in\Couplings(\al,\be)$, disintegrate $\pi$ with respect to
  each marginal and define
  \begin{align}
    T_\pi:L^2(\be)&\longrightarrow L^2(\al),
    &T_\pi k(x)&:=\E_\pi[k(Y)\mid X=x],
    \label{eq:Tpi-definition}
    \\
    T_\pi^\ast:L^2(\al)&\longrightarrow L^2(\be),
    &T_\pi^\ast h(y)&:=\E_\pi[h(X)\mid Y=y].
    \label{eq:Tpi-adjoint-definition}
  \end{align}
  Its centered maximal-correlation coefficient is
  \begin{align}
    \sigma(\pi)
    &:=
    \opnorm{T_\pi:L_0^2(\be)\to L_0^2(\al)}
    \label{eq:sigma-pi-definition}
    \\
    &=
    \sup_{\substack{h\in L_0^2(\al),\ k\in L_0^2(\be)\\
                     \norm{h}_\al=\norm{k}_\be=1}}
    \int_{\Xx\times\Yy}h(x)k(y)\,\dd\pi(x,y).
    \notag
  \end{align}
  For the entropic optimizer \eqref{eq:optimal-gibbs-plan}, set
  \begin{equation}
    \label{eq:optimal-coupling-operator}
    T_\eps:=T_{\pi_\eps^\star},
    \qquad
    \sigma_\eps:=\sigma(\pi_\eps^\star).
  \end{equation}
\end{definition}

The adjoint notation is literal:
$\ip{h}{T_\pi k}_{\al}=\ip{T_\pi^\ast h}{k}_{\be}$.
Conditional Jensen gives $\opnorm{T_\pi}\leq1$, while
$T_\pi\one=\one$ and $T_\pi^\ast\one=\one$. Hence $T_\pi$ and
$T_\pi^\ast$ preserve centered functions and $\sigma(\pi)\in[0,1]$.
The spectral gap of $T_\eps^\ast T_\eps$ on $L_0^2(\be)$ is
$1-\sigma_\eps^2$.

\begin{proposition}[Dual Hessian at arbitrary and optimal potentials]
  \label{prop:full-dual-hessian}
  For arbitrary potentials $(f,g)$ and directions $(h,k),(h',k')$ for
  which the expressions are finite,
  \begin{equation}
    \label{eq:general-dual-second-variation}
    -\eps D^2\Dd_\eps(f,g)[(h,k),(h',k')]
    =
    \int_{\Xx\times\Yy}
    (h(x)+k(y))(h'(x)+k'(y))\,\dd\pi_{f,g}(x,y).
  \end{equation}
  At an optimal pair $(f^\star,g^\star)$, the positive curvature
  operator on $L^2(\al)\oplus L^2(\be)$ is
  \begin{equation}
    \label{eq:full-dual-hessian}
    H_{\mathrm{dual}}
    =
    -\nabla^2\Dd_\eps(f^\star,g^\star)
    =
    \frac1\eps
    \begin{pmatrix}
      I&T_{\pi_\eps^\star}\\
      T_{\pi_\eps^\star}^\ast&I
    \end{pmatrix}.
  \end{equation}
  The Hessian descends to the quotient
  $\mathcal H_{\mathrm q}$ in \eqref{eq:dual-quotient-space}, where
  \begin{equation}
    \label{eq:full-dual-quotient-curvature}
    \frac{1-\sigma_\eps}{\eps}I
    \preceq H_{\mathrm{dual}}
    \preceq
    \frac{2}{\eps}I.
  \end{equation}
  Consequently, if $\sigma_\eps<1$, its spectral condition number is
  \begin{equation}
    \label{eq:full-dual-condition}
    \kappa_{\mathrm{dual}}
    =
    \frac{2}{1-\sigma_\eps}.
  \end{equation}
  If $\sigma_\eps=1$, the quotient Hessian is not coercive and this
  condition number is infinite.
\end{proposition}

\begin{proof}
Twice differentiating the exponential term in $-\Dd_\eps$ gives
\eqref{eq:general-dual-second-variation}. At the optimum,
$\pi_{f^\star,g^\star}=\pi_\eps^\star$ has marginals $(\al,\be)$.
The two diagonal terms in \eqref{eq:general-dual-second-variation} are
therefore the $L^2(\al)$ and $L^2(\be)$ inner products, while the cross
term is $\ip{h}{T_{\pi_\eps^\star}k}_\al$. This proves
\eqref{eq:full-dual-hessian}. On a centered singular pair
$T_\eps v=s u$, $T_\eps^\ast u=s v$, the two vectors $(u,\pm v)$ have
eigenvalues $(1\pm s)/\eps$. A representative orthogonal to the gauge
decomposes uniquely as
\[
  (h,k)=(h_0,k_0)+m(\one,\one),
  \qquad
  h_0\in L_0^2(\al),\quad k_0\in L_0^2(\be).
\]
The two summands are orthogonal and invariant under the Hessian. The
first has spectral interval
$[(1-\sigma_\eps)/\eps,(1+\sigma_\eps)/\eps]$, while the second has
eigenvalue $2/\eps$. Since $1+\sigma_\eps\leq2$, this proves
\eqref{eq:full-dual-quotient-curvature} and
\eqref{eq:full-dual-condition}.
\end{proof}

Thus optimality is not needed for the second-variation identity
\eqref{eq:general-dual-second-variation}. It is needed for the clean
conditional-expectation form \eqref{eq:full-dual-hessian}: away from
stationarity, $\pi_{f,g}$ is generally not a coupling, the diagonal
blocks are multiplication by its marginal densities relative to
$\al$ and $\be$, and the off-diagonal block is an unnormalized integral
operator rather than $T_\pi$ from Definition~\ref{def:coupling-operator}.

For discrete marginals $\al=\sum_i \mathrm a_i\delta_{x_i}$ and
$\be=\sum_j \mathrm b_j\delta_{y_j}$, with optimal coupling matrix
$\mathrm P_\eps^\star$, the same operators read
\[
  T_{\pi_\eps^\star}k
  =\mathrm D_{\mathrm a}^{-1}\mathrm P_\eps^\star k,
  \qquad
  T_{\pi_\eps^\star}^\ast h
  =\mathrm D_{\mathrm b}^{-1}
  (\mathrm P_\eps^\star)^{\mathsf T}h,
\]
where $\mathrm D_{\mathrm a}=\operatorname{diag}(\mathrm a)$ and
$\mathrm D_{\mathrm b}=\operatorname{diag}(\mathrm b)$. These formulas
are adjoint in the $\mathrm a$- and $\mathrm b$-weighted inner products,
not in the unweighted Euclidean ones.

\section{Sinkhorn Jacobian and local rate}

A full Sinkhorn cycle is a nonlinear map on one potential. Its
linearization at the fixed point is exactly the two-step conditional
expectation $T_\eps^\ast T_\eps$, so the local convergence question reduces to a
Markov spectral gap.

\paragraph{Soft transforms and the cycle map.}
Exact block maximization of $\Dd_\eps$ is written with the continuous
soft $c$-transforms
\begin{align}
  g^{\bar c,\eps}(x)
  &:=-\eps\log\int_{\Yy}
  \exp\left(\frac{g(y)-c(x,y)}{\eps}\right)\dd\be(y),
  \label{eq:soft-transform-g}
  \\
  f^{c,\eps}(y)
  &:=-\eps\log\int_{\Xx}
  \exp\left(\frac{f(x)-c(x,y)}{\eps}\right)\dd\al(x).
  \label{eq:soft-transform-f}
\end{align}
One complete cycle on the target potential is therefore
\begin{equation}
  \label{eq:sinkhorn-cycle}
  \Ss_\eps(g):=(g^{\bar c,\eps})^{c,\eps},
  \qquad
  g^{(\ell+1)}=\Ss_\eps(g^{(\ell)}).
\end{equation}
The opposite ordering acts on $L^2(\al)$ and has the same nonzero
spectrum.

\begin{proposition}[Jacobian and exact local Sinkhorn rate]
  \label{prop:sinkhorn-local-rate}
  Assume that the two soft transforms are twice differentiable near
  the optimum. Then
  \begin{equation}
    \label{eq:coordinate-jacobians}
    D[g\mapsto g^{\bar c,\eps}](g^\star)=-T_\eps,
    \qquad
    D[f\mapsto f^{c,\eps}](f^\star)=-T_\eps^\ast,
    \qquad
    D\Ss_\eps(g^\star)=T_\eps^\ast T_\eps.
  \end{equation}
  Consequently, for the centered error
  $e^{(\ell)}=\operatorname{cent}_{\be}(g^{(\ell)}-g^\star)$ of any Sinkhorn sequence
  converging to $g^\star$ modulo constants,
  \begin{equation}
    \label{eq:local-error-recursion}
    e^{(\ell+1)}
    =
    T_\eps^\ast T_\eps e^{(\ell)}+o(\norm{e^{(\ell)}}_\be),
    \qquad
    \limsup_{\ell\to\infty}
    \frac{\norm{e^{(\ell+1)}}_\be}
         {\norm{e^{(\ell)}}_\be}
    \leq \sigma_\eps^2.
  \end{equation}
  If $\sigma_\eps^2$ is an isolated eigenvalue and the asymptotic error
  has a nonzero projection onto its eigenspace, the ratio converges to
  $\sigma_\eps^2$.
\end{proposition}

\begin{proof}
Differentiating the log-partition maps
\eqref{eq:soft-transform-g}--\eqref{eq:soft-transform-f} at the optimal
Gibbs coupling gives the first two identities in
\eqref{eq:coordinate-jacobians}; the
chain rule gives the third. Taylor expansion proves
\eqref{eq:local-error-recursion}.
\end{proof}

The reverse cycle $f\mapsto(f^{c,\eps})^{\bar c,\eps}$ has Jacobian
$T_\eps T_\eps^\ast$ on
$L_0^2(\al)$. Since $T_\eps^\ast T_\eps$ and
$T_\eps T_\eps^\ast$ have the same nonzero
spectrum, either orientation has the same local rate. In finite
spaces, this recovers the classical second-singular-value formula for
matrix scaling \cite{Knight2008,LehmannEtAl2022}.

\paragraph{When is the rate strictly smaller than one?}
Positivity of the Gibbs kernel is enough in finite dimension but not on
an arbitrary infinite space: centered approximate modes may retain
correlation arbitrarily close to one. In the finite or compact
continuous setting, strictness follows directly from the same
Birkhoff--Hilbert contraction theorem that gives global linear
convergence of Sinkhorn. The only additional step is to differentiate
that global estimate at the fixed point.

Let $k_\eps(x,y):=e^{-c(x,y)/\eps}$. For compact supports and continuous
$c$, define the four-point oscillation
\begin{equation}
  \label{eq:four-point-cost-oscillation}
  \Delta_c
  :=
  \sup_{x,x',y,y'}
  \bigl[c(x,y')+c(x',y)-c(x,y)-c(x',y')\bigr],
\end{equation}
where the supremum is over the product of the two supports.
This quantity is invariant under $c(x,y)\mapsto
c(x,y)+\xi(x)+\zeta(y)$ and is exactly the projective diameter of the
Gibbs operator, up to temperature:
\begin{equation}
  \label{eq:gibbs-projective-diameter}
  \sup_{x,x',y,y'}
  \log
  \frac{k_\eps(x,y)k_\eps(x',y')}
       {k_\eps(x,y')k_\eps(x',y)}
  =
  \frac{\Delta_c}{\eps}.
\end{equation}

\begin{proposition}[Birkhoff--Hilbert bound on the local rate]
  \label{prop:strict-gap}
  Assume either finite spaces with a positive Gibbs kernel, or compact
  supports with continuous finite $c$. Set
  \begin{equation}
    \label{eq:birkhoff-local-bound}
    \Lambda_\eps
    :=
    \tanh\left(\frac{\Delta_c}{4\eps}\right)<1.
  \end{equation}
  Each Gibbs half-step is $\Lambda_\eps$-contractive in Hilbert's
  projective metric. Consequently, a complete Sinkhorn cycle is globally
  $\Lambda_\eps^2$-contractive in Hilbert's projective metric, and its
  local $L^2$ factor satisfies
  \begin{equation}
    \label{eq:birkhoff-controls-sigma}
    \sigma_\eps^2\leq\Lambda_\eps^2,
    \qquad\text{hence}\qquad
    \sigma_\eps\leq\Lambda_\eps<1.
  \end{equation}
\end{proposition}

\begin{proof}
Equation~\eqref{eq:gibbs-projective-diameter} and the Birkhoff--Hopf
theorem give the half-step estimate
\eqref{eq:birkhoff-local-bound}; multiplication and inversion in the
scaling updates are projective isometries, so a full cycle has factor
$\Lambda_\eps^2$
\cite{FranklinLorenz1989}. In logarithmic coordinates, Hilbert's metric
is the variation seminorm. Differentiating the full-cycle contraction
at its fixed point and using Proposition~\ref{prop:sinkhorn-local-rate}
therefore gives
\[
  \operatorname{osc}(T_\eps^\ast T_\eps k)
  \leq
  \Lambda_\eps^2\operatorname{osc}(k).
\]
In finite dimension this applies directly to a top centered
eigenvector of $T_\eps^\ast T_\eps$. In the compact continuous case,
the conditional operator is compact and that eigenvector is continuous.
Its eigenvalue is $\sigma_\eps^2$, which proves
\eqref{eq:birkhoff-controls-sigma}; the case $\sigma_\eps=0$ is
immediate.
\end{proof}

A cost-oscillation estimate modulo separable terms is only a coarser
way to bound the same Birkhoff coefficient: for every $\xi,\zeta$,
\[
  \Delta_c
  \leq
  2\,\operatorname{ess\,osc}_{\al\otimes\be}
  \bigl(c(x,y)-\xi(x)-\zeta(y)\bigr).
\]
Thus it introduces no distinct contraction mechanism and need not be
proved separately. On unbounded spaces $\Delta_c$ is generally
infinite, so the universal Birkhoff estimate becomes vacuous. Cost
growth must then be balanced against marginal tails. Conforti, Durmus,
and Greco prove quantitative contraction for quadratic costs under
asymptotically positive log-concavity profiles
\cite{ConfortiDurmusGreco2026}; Eckstein obtains weighted Hilbert
contraction under compatible cost-growth and tail assumptions
\cite{Eckstein2025}. These are useful input-level criteria, but their
norms should not be identified automatically with the exact $L^2$ gap
$1-\sigma_\eps^2$.

\paragraph{Why a linear small-$\eps$ gap is natural.}
The conditional-variance identity becomes useful specifically for
understanding the small-regularization regime. The law of total
variance gives
\begin{equation}
  \label{eq:conditional-poincare-gap}
  1-\sigma_\eps^2
  =
  \inf_{\substack{k\in L_0^2(\be)\\\norm{k}_\be=1}}
  \E_{\pi_\eps^\star}\left[\Var(k(Y)\mid X)\right].
\end{equation}
Indeed, the right-hand side is the smallest Rayleigh quotient of
$I-T_\eps^\ast T_\eps$ on $L_0^2(\be)$. In a smooth connected Monge
regime, the conditional law of $Y$ given $X=x$ has width
$O(\sqrt\eps)$ around the hard transport map. Consequently, the
conditional variance in \eqref{eq:conditional-poincare-gap} is
typically of order $\eps$, leading to
$\sigma_\eps^2=1-C\eps+o(\eps)$.

A precise conditional statement is as follows. Assume that $\Xx$ and $\Yy$
are smooth connected $d$-dimensional manifolds, that the marginals have
positive smooth densities, and that hard transport is induced by a
smooth diffeomorphism $\mathsf M:\Xx\to\Yy$. Let $g_0^\star$ be the hard
target potential and set
\begin{equation}
  \label{eq:monge-transverse-hessian}
  H(x)
  :=
  D_{yy}^2\bigl(c(x,\cdot)-g_0^\star\bigr)(\mathsf M(x)),
  \qquad
  A(y):=H(\mathsf M^{-1}(y))^{-1}.
\end{equation}
Suppose that $H$ is uniformly positive and that the uniform Laplace
expansion of the entropic conditionals permits passage to the first
nonzero Rayleigh quotient, for instance through Mosco convergence of
the rescaled forms and compact convergence of their resolvents. Define
\begin{equation}
  \label{eq:limiting-poincare-eigenvalue}
  \lambda_1(A,\be)
  :=
  \inf_{\substack{\int k\dd\be=0\\\int k^2\dd\be=1}}
  \int_{\Yy}\ip{A(y)\nabla k(y)}{\nabla k(y)}\dd\be(y).
\end{equation}
Uniform ellipticity and a Poincar\'e inequality for $\be$ ensure that
$\lambda_1(A,\be)>0$.

\begin{proposition}[Smooth small-regularization law]
  \label{prop:smooth-small-epsilon}
  Under the preceding assumptions,
  \begin{equation}
    \label{eq:smooth-small-epsilon-rate}
    \sigma_\eps^2
    =
    1-\eps\lambda_1(A,\be)+o(\eps),
    \qquad
    \sigma_\eps
    =
    1-\frac\eps2\lambda_1(A,\be)+o(\eps).
  \end{equation}
\end{proposition}

\begin{proof}
For the optimal coupling $\pi_\eps^\star$, the uniform Laplace expansion
of $Y\mid X=x$ gives
$\Cov(Y\mid X=x)=\eps H(x)^{-1}+o(\eps)$. Hence, for smooth $k$,
\[
  \frac1\eps
  \E_{\pi_\eps^\star}[\Var(k(Y)\mid X)]
  \longrightarrow
  \int_{\Xx}
  \ip{H(x)^{-1}\nabla k(\mathsf M(x))}
     {\nabla k(\mathsf M(x))}\dd\al(x).
\]
Since $\mathsf M_\#\al=\be$, the right-hand side is the form in
\eqref{eq:limiting-poincare-eigenvalue}. Spectral convergence permits
passage to the first Rayleigh quotient in
\eqref{eq:conditional-poincare-gap}, giving the expansion of
$\sigma_\eps^2$. The expansion of $\sigma_\eps$ follows by taking the
square root.
\end{proof}

\begin{example}[Exact Gaussian benchmark]
  \label{ex:gaussian-benchmark}
  For $\al=\be=\Gaussian(0,1)$ and
  $c(x,y)=\frac12(x-y)^2$, the optimal coupling is Gaussian with
  correlation
  \begin{equation}
    \label{eq:gaussian-correlation}
    r_\eps
    =
    \frac{\sqrt{\eps^2+4}-\eps}{2}.
  \end{equation}
  Hermite polynomials diagonalize $T_\eps$, so $\sigma_\eps=r_\eps$.
  Thus $\sigma_\eps=1-\eps/2+O(\eps^2)$ and
  $\sigma_\eps^2=1-\eps+O(\eps^2)$. This is
  Proposition~\ref{prop:smooth-small-epsilon} with
  $A=1$ and $\lambda_1(A,\be)=1$.
\end{example}

\paragraph{Local $L^2$ versus global Hilbert contraction.}
For $\Delta_c>0$,
\begin{equation}
  \label{eq:hilbert-small-epsilon}
  1-\Lambda_\eps^2
  \sim
  4e^{-\Delta_c/(2\eps)}
  \qquad(\eps\downarrow0).
\end{equation}
This is exponentially more pessimistic than the smooth local law
$1-\sigma_\eps^2\sim\eps\lambda_1(A,\be)$. There is no contradiction.
Hilbert's projective metric gives a global guarantee over the full positive cone and uses
no smoothness of the iterates. The $L^2$ result is local to the optimal
coupling and exploits the regular diffusive structure of the dominant
modes. Chizat, Delalande, and Va\v{s}kevi\v{c}ius similarly obtain
polynomial rather than universal exponential dependences under
continuous regularity and log-concavity assumptions
\cite{ChizatDelalandeVaskevicius2026}.

\section{Blockwise over-relaxation}
\label{sec:acceleration}

Sinkhorn alternates two exact block maximizations. A direct acceleration
is therefore to extrapolate each block separately toward its soft
$c$-transform. This is the nonlinear successive-over-relaxation (SOR)
scheme studied in
\cite{ThibaultEtAl2021,LehmannEtAl2022}.

\paragraph{Scalar extrapolation of the two block updates.}
Starting from $(f^{(\ell)},g^{(\ell)})$, and for a relaxation parameter
$\omega>0$, set
\begin{equation}
  \label{eq:block-sor}
  \begin{aligned}
    f^{(\ell+1)}
    &=
    (1-\omega)f^{(\ell)}
    +\omega\bigl(g^{(\ell)}\bigr)^{\bar c,\eps},
    \\
    g^{(\ell+1)}
    &=
    (1-\omega)g^{(\ell)}
    +\omega\bigl(f^{(\ell+1)}\bigr)^{c,\eps}.
  \end{aligned}
\end{equation}
The choice $\omega=1$ recovers ordinary Sinkhorn; $0<\omega<1$
under-relaxes both updates, while $1<\omega<2$ extrapolates beyond each
block maximizer. The following result gives the exact local spectral
factor.

\begin{proposition}[Local rate and optimal block relaxation]
  \label{prop:optimal-overrelaxation}
  Assume the differentiability hypotheses of
  Proposition~\ref{prop:sinkhorn-local-rate} and $\sigma_\eps<1$.
  For $0<\omega<2$, define
  \begin{equation}
    \label{eq:sor-threshold}
    \omega_\star
    :=
    \frac{2}{1+\sqrt{1-\sigma_\eps^2}},
    \qquad
    \tau_\omega
    :=
    2(1-\omega)+\omega^2\sigma_\eps^2.
  \end{equation}
  The local asymptotic factor of one complete iteration
  \eqref{eq:block-sor}, modulo the additive gauge, is
  \begin{equation}
    \label{eq:sor-rate}
    r_\omega
    =
    \begin{cases}
      \dfrac{\tau_\omega+
      \sqrt{\tau_\omega^2-4(1-\omega)^2}}{2},
      & 0<\omega\leq\omega_\star,
      \\[3mm]
      \omega-1,
      & \omega_\star\leq\omega<2.
    \end{cases}
  \end{equation}
  It is minimized by $\omega=\omega_\star$, with optimal factor
  \begin{equation}
    \label{eq:optimal-sor-rate}
    r_\star
    =
    \omega_\star-1
    =
    \frac{1-\sqrt{1-\sigma_\eps^2}}
         {1+\sqrt{1-\sigma_\eps^2}}.
  \end{equation}
  In particular, $\omega=1$ gives $r_1=\sigma_\eps^2$, the ordinary
  Sinkhorn rate.
\end{proposition}

\begin{proof}
On centered functions, use the polar decomposition of $T_\eps$ and the
spectral theorem for $(T_\eps^\ast T_\eps)^{1/2}$. For each spectral
singular value $s\in[0,\sigma_\eps]$, linearizing
\eqref{eq:block-sor} with \eqref{eq:coordinate-jacobians} reduces the
error evolution to
\begin{equation}
  \label{eq:sor-mode-matrix}
  \mathsf M_\omega(s)
  :=
  \begin{pmatrix}
    1-\omega & -\omega s\\
    -\omega s(1-\omega) & 1-\omega+\omega^2s^2
  \end{pmatrix}.
\end{equation}
Its characteristic polynomial is
\[
  z^2-\bigl(2(1-\omega)+\omega^2s^2\bigr)z
  +(1-\omega)^2.
\]
For $0<\omega\leq\omega_\star$, the largest root modulus is attained at
$s=\sigma_\eps$ and equals the first branch of
\eqref{eq:sor-rate}. The threshold $\omega_\star$ is the smaller
solution of
$\omega^2\sigma_\eps^2=4(\omega-1)$. For
$\omega_\star\leq\omega<2$, the two roots are conjugate for every
$0<s\leq\sigma_\eps$ and have modulus $\omega-1$; at $s=0$ the
repeated root has the same modulus. The nongauge constant mode has factor
$(1-\omega)^2$ and never dominates. Thus \eqref{eq:sor-rate} is the
spectral radius of the linearized cycle. Its first branch decreases up
to $\omega_\star$, whereas its second branch increases, proving
\eqref{eq:optimal-sor-rate}.
\end{proof}

When $1-\sigma_\eps^2$ is small, the optimal factor satisfies
$r_\star=1-2\sqrt{1-\sigma_\eps^2}
+O(1-\sigma_\eps^2)$, instead of the ordinary rate
$\sigma_\eps^2$. This is a local prescription: aggressive
over-relaxation need not be globally convergent without the additional
projective bounds or safeguards developed in
\cite{ThibaultEtAl2021,LehmannEtAl2022}.

\section{Semi-dual variable projection}
\label{sec:semidual}

Eliminating one potential exactly converts the same operator
$T_\eps^\ast T_\eps$ from an iteration Jacobian into a Hessian correction. This
is the variable-projection mechanism: the poorly conditioned block
coupling is removed by a Schur complement rather than handled by
alternating block maximization.

This elimination is a special case of the variable-projection method
of Golub and Pereyra
\cite{GolubPereyra1973,GolubPereyra2003}. In general, if one block is
optimized exactly (or minimized for the negative objective) and the
positive Hessian at a stationary point is
$H=\left(\begin{smallmatrix}H_{11}&H_{12}\\
H_{21}&H_{22}\end{smallmatrix}\right)$ with $H_{11}\succ0$, then the
reduced Hessian is
\begin{equation}
  \label{eq:general-varpro-schur}
  H_{\mathrm{red}}
  =
  H_{22}-H_{21}H_{11}^{-1}H_{12}.
\end{equation}
Its variational characterization gives, with extremal spectral values
in infinite dimension,
$\lambda_{\min}(H)\leq\lambda_{\min}(H_{\mathrm{red}})$ and
$\lambda_{\max}(H_{\mathrm{red}})\leq\lambda_{\max}(H)$, so exact
projection cannot worsen the local spectral condition number after the
gauge has been removed. General variable-projection theory does not,
however, quantify this gain without additional information on the
blocks. Here both \emph{diagonal} blocks of the normalized full-dual
Hessian are the identity, while the off-diagonal blocks are
$T_\eps$ and $T_\eps^\ast$. This special structure makes the reduced
operator $I-T_\eps^\ast T_\eps$ and its conditioning explicitly
computable from $\sigma_\eps$.

\paragraph{Reduced objective and Hessian.}
For each $g$, the exact source-potential maximizer is
$g^{\bar c,\eps}$ from \eqref{eq:soft-transform-g}. The resulting
concave semi-dual is
\begin{equation}
  \label{eq:semidual-definition}
  \Ee_\eps(g)
  :=
  \Dd_\eps(g^{\bar c,\eps},g)
  =
  \int_{\Xx}g^{\bar c,\eps}\dd\al
  +
  \int_{\Yy}g\dd\be.
\end{equation}
Let
\begin{equation}
  \label{eq:semidual-coupling}
  \dd\pi_g(x,y)
  :=
  \exp\left(
    \frac{g^{\bar c,\eps}(x)+g(y)-c(x,y)}{\eps}
  \right)
  \dd\al(x)\dd\be(y),
\end{equation}
and denote its second marginal by $\be_g$. Its first marginal is
exactly $\al$.

\begin{proposition}[Semi-dual derivatives and Schur complement]
  \label{prop:semidual-varpro}
  For bounded directions $k,k'$ on $\Yy$,
  \begin{align}
    D\Ee_\eps(g)[k]
    &=
    \int_{\Yy} k\,\dd(\be-\be_g),
    \label{eq:semidual-gradient}
    \\
    D^2\Ee_\eps(g)[k,k']
    &=
    -\frac1\eps
    \int_{\Xx}
    \Cov_{\pi_g(\cdot\mid x)}(k(Y),k'(Y))\dd\al(x).
    \label{eq:semidual-hessian-general}
  \end{align}
  At $g=g^\star$, the positive Riesz curvature operator on $L^2(\be)$ is
  \begin{equation}
    \label{eq:semidual-hessian}
    H_{\mathrm{sd}}
    :=
    -\nabla^2\Ee_\eps(g^\star)
    =
    \frac1\eps(I-T_\eps^\ast T_\eps).
  \end{equation}
  The identity extends from bounded directions to $L^2(\be)$ by
  density.
  It is exactly the Schur complement of the $f$ block in
  \eqref{eq:full-dual-hessian}.
\end{proposition}

\begin{proof}
The envelope theorem gives \eqref{eq:semidual-gradient}. Differentiating
the conditional Gibbs law in \eqref{eq:semidual-coupling} gives the
conditional covariance formula \eqref{eq:semidual-hessian-general}.
At the optimum,
\[
  \int_{\Xx}\Cov(k,k'\mid X)\dd\al
  =
  \ip{k}{k'}_\be-\ip{T_\eps k}{T_\eps k'}_\al
  =
  \ip{k}{(I-T_\eps^\ast T_\eps)k'}_\be,
\]
which proves \eqref{eq:semidual-hessian}. Finally, the Schur complement
of $I$ in
$\left(\begin{smallmatrix}I&T_\eps\\T_\eps^\ast&I\end{smallmatrix}\right)$ is
$I-T_\eps^\ast T_\eps$, which is
\eqref{eq:general-varpro-schur} in the present geometry.
\end{proof}

\paragraph{Conditioning and first-order rates.}
The spectral gap that controls Sinkhorn is exactly the strong concavity
of $\Ee_\eps$, or equivalently the strong convexity of $-\Ee_\eps$.

\begin{proposition}[Semi-dual conditioning and gradient-ascent rate]
  \label{prop:semidual-conditioning}
  On $L_0^2(\be)$,
  \begin{equation}
    \label{eq:semidual-curvature}
    \frac{1-\sigma_\eps^2}{\eps} I
    \preceq
    H_{\mathrm{sd}}
    \preceq
    \frac1\eps I,
    \qquad
    \kappa_{\mathrm{sd}}
    \leq
    \frac1{1-\sigma_\eps^2}.
  \end{equation}
  Gradient ascent on the local quadratic model of $\Ee_\eps$, with
  optimal constant step $2\eps/(2-\sigma_\eps^2)$, has factor
  \begin{equation}
    \label{eq:semidual-gradient-rate}
    r_{\mathrm{GD}}
    \leq
    \frac{\sigma_\eps^2}{2-\sigma_\eps^2}.
  \end{equation}
\end{proposition}

\begin{proof}
The spectrum of $T_\eps^\ast T_\eps$ on centered functions lies in
$[0,\sigma_\eps^2]$, so
the spectrum of \eqref{eq:semidual-hessian} lies in
$[(1-\sigma_\eps^2)/\eps,1/\eps]$. This proves the curvature and
condition-number bounds. The standard minimax polynomials on an interval give
\eqref{eq:semidual-gradient-rate}.
\end{proof}

Compared with the quotient full-dual condition number
\eqref{eq:full-dual-condition}, the displayed semi-dual upper bound is
smaller by the factor $2(1+\sigma_\eps)$. This is the quantitative gain
furnished here by variable projection.

\paragraph{Quasi-Newton methods.}
In a finite discretization, $-\Ee_\eps$ is a smooth unconstrained
convex problem after fixing its additive gauge. BFGS attempts to learn
the inverse of $I-T_\eps^\ast T_\eps$ from gradient differences, precisely the
operator responsible for the slow modes. Under the standard local
smoothness and positive-definiteness assumptions, full BFGS is locally
superlinear; limited-memory BFGS trades that guarantee for lower memory
and often strong practical performance \cite{NocedalWright2006}.

This strategy has a substantial OT history. Cuturi and Peyr\'e used
smooth dual and semi-dual formulations with L-BFGS for variational
Wasserstein problems and developed the general regularized semi-dual
viewpoint \cite{CuturiPeyre2016,CuturiPeyre2018}; stochastic
first-order optimization of semi-duals was developed by Genevay et al.
\cite{GenevayEtAl2016}. Khamlich, Romor, and
Rozza use L-BFGS as the optimization core of a modern
entropy-regularized semi-discrete solver, together with truncation,
multilevel continuation, and fast geometric queries
\cite{KhamlichRomorRozza2026}. Qiu, Yin, and Wang also report and analyze
L-BFGS on the fully discrete entropic dual as an alternative to long
Sinkhorn runs \cite{QiuYinWang2023}.

These methods do not make the small-$\eps$ conditioning disappear: the
inverse Hessian still grows like $1/(1-\sigma_\eps^2)$. Their advantage
is that curvature information is learned and reused. Each quasi-Newton
step is more expensive than one scaling update, so comparisons should
count kernel evaluations or matrix-vector products rather than
iterations. The semi-dual is especially attractive when one block can
be eliminated analytically or integrated accurately, as in
semi-discrete transport.

\begin{remark}[Blockwise over-relaxation versus semi-dual gradient ascent]
  \label{rem:sinkhorn-versus-semdual-gradient}
  Even ordinary Sinkhorn and semi-dual gradient ascent are not identical
  globally. For the coupling
  $\pi_g$ in \eqref{eq:semidual-coupling}, let
  $r_g:=\dd\be_g/\dd\be$. The scaling identities give
  \begin{equation}
    \label{eq:sinkhorn-gradient-global-comparison}
    r_g
    =
    \exp\left(\frac{g-\Ss_\eps(g)}{\eps}\right),
    \qquad
    \nabla\Ee_\eps(g)=1-r_g,
    \qquad
    \Ss_\eps(g)-g=-\eps\log r_g.
  \end{equation}
  Thus a Sinkhorn cycle uses the logarithmic marginal defect, whereas
  semi-dual gradient ascent uses the additive defect. They agree only
  to first order near $r_g=1$, where
  \begin{equation}
    \label{eq:sinkhorn-gradient-local-comparison}
    \Ss_\eps(g)-g
    =
    \eps\nabla\Ee_\eps(g)
    +o(\norm{g-g^\star}).
  \end{equation}
  Thus ordinary Sinkhorn is locally semi-dual gradient ascent with step
  $\eps$. The blockwise scheme \eqref{eq:block-sor} is different: when
  $\omega\neq1$, the source potential is only partially updated and
  remains an independent state. It is therefore neither an outer
  extrapolation of $\Ss_\eps$ nor a semi-dual gradient step with step
  $\eps\omega$; its local Jacobian is the two-block SOR matrix
  \eqref{eq:sor-mode-matrix}. Accordingly, its optimal factor
  \eqref{eq:optimal-sor-rate} has a square-root dependence on
  $1-\sigma_\eps^2$, whereas optimal fixed-step semi-dual gradient
  ascent has the factor \eqref{eq:semidual-gradient-rate}. This is a
  local spectral comparison, not a global equivalence of the nonlinear
  algorithms.
\end{remark}

\paragraph{Operator summary.}

All local statements in the note follow from one contraction:
\begin{center}
\begin{tabular}{@{}lll@{}}
\toprule
Object & Operator at the optimum & Consequence \\
\midrule
Full dual Hessian
& $\eps^{-1}\left(\begin{smallmatrix}I&T_\eps\\T_\eps^\ast&I\end{smallmatrix}\right)$
& coupled conditioning \\
Sinkhorn full-cycle Jacobian
& $T_\eps^\ast T_\eps$
& local factor $\sigma_\eps^2$ \\
Semi-dual Hessian
& $\eps^{-1}(I-T_\eps^\ast T_\eps)$
& curvature $(1-\sigma_\eps^2)/\eps$ \\
Blockwise over-relaxation
& $\mathsf M_\omega(s)$ from \eqref{eq:sor-mode-matrix}
& optimal factor $r_\star$ in \eqref{eq:optimal-sor-rate} \\
\bottomrule
\end{tabular}
\end{center}
Thus the slow convergence of Sinkhorn and the ill-conditioning of the
semi-dual are complementary descriptions of the same near-deterministic
correlation modes. Variable projection removes one coupled block;
blockwise over-relaxation modifies the alternating splitting; and
quasi-Newton methods approximate the inverse semi-dual curvature.

\begin{thebibliography}{99}

\bibitem{ChizatDelalandeVaskevicius2026}
L. Chizat, A. Delalande, and T. Va\v{s}kevi\v{c}ius,
\newblock Sharper exponential convergence rates for Sinkhorn's
algorithm in continuous settings,
\newblock \emph{Mathematical Programming}, 215:809--858, 2026.
\newblock DOI: 10.1007/s10107-025-02242-z.

\bibitem{ConfortiDurmusGreco2026}
G. Conforti, A. Durmus, and G. Greco,
\newblock Quantitative contraction rates for Sinkhorn's algorithm:
beyond bounded costs and compact marginals,
\newblock \emph{Annals of Applied Probability}, to appear, 2026.
\newblock arXiv:2304.04451.

\bibitem{Cuturi2013}
M. Cuturi,
\newblock Sinkhorn distances: Lightspeed computation of optimal
transport,
\newblock in \emph{Advances in Neural Information Processing Systems
26}, pages 2292--2300, 2013.

\bibitem{CuturiPeyre2016}
M. Cuturi and G. Peyr\'e,
\newblock A smoothed dual approach for variational Wasserstein
problems,
\newblock \emph{SIAM Journal on Imaging Sciences},
9(1):320--343, 2016.
\newblock DOI: 10.1137/15M1032600.

\bibitem{CuturiPeyre2018}
M. Cuturi and G. Peyr\'e,
\newblock Semidual regularized optimal transport,
\newblock \emph{SIAM Review}, 60(4):941--965, 2018.
\newblock DOI: 10.1137/18M1208654.

\bibitem{Eckstein2025}
S. Eckstein,
\newblock Hilbert's projective metric for functions of bounded growth
and exponential convergence of Sinkhorn's algorithm,
\newblock \emph{Probability Theory and Related Fields},
193:585--621, 2025.
\newblock DOI: 10.1007/s00440-025-01366-9.

\bibitem{FranklinLorenz1989}
J. Franklin and J. Lorenz,
\newblock On the scaling of multidimensional matrices,
\newblock \emph{Linear Algebra and its Applications},
114--115:717--735, 1989.

\bibitem{GolubPereyra1973}
G. H. Golub and V. Pereyra,
\newblock The differentiation of pseudo-inverses and nonlinear least
squares problems whose variables separate,
\newblock \emph{SIAM Journal on Numerical Analysis},
10(2):413--432, 1973.
\newblock DOI: 10.1137/0710036.

\bibitem{GolubPereyra2003}
G. H. Golub and V. Pereyra,
\newblock Separable nonlinear least squares: The variable projection
method and its applications,
\newblock \emph{Inverse Problems}, 19(2):R1--R26, 2003.

\bibitem{GenevayEtAl2016}
A. Genevay, M. Cuturi, G. Peyr\'e, and F. Bach,
\newblock Stochastic optimization for large-scale optimal transport,
\newblock in \emph{Advances in Neural Information Processing Systems
29}, 2016.

\bibitem{KhamlichRomorRozza2026}
M. Khamlich, F. Romor, and G. Rozza,
\newblock Efficient numerical strategies for entropy-regularized
semi-discrete optimal transport,
\newblock \emph{Computer Methods in Applied Mechanics and Engineering},
453:118821, 2026.
\newblock DOI: 10.1016/j.cma.2026.118821.

\bibitem{Knight2008}
P. A. Knight,
\newblock The Sinkhorn--Knopp algorithm: Convergence and applications,
\newblock \emph{SIAM Journal on Matrix Analysis and Applications},
30(1):261--275, 2008.
\newblock DOI: 10.1137/060659624.

\bibitem{LehmannEtAl2022}
T. Lehmann, M.-K. von Renesse, A. Sambale, and A. Uschmajew,
\newblock A note on overrelaxation in the Sinkhorn algorithm,
\newblock \emph{Optimization Letters}, 16(8):2209--2220, 2022.
\newblock DOI: 10.1007/s11590-021-01830-0.

\bibitem{NocedalWright2006}
J. Nocedal and S. J. Wright,
\newblock \emph{Numerical Optimization}, second edition,
\newblock Springer, 2006.

\bibitem{QiuYinWang2023}
Y. Qiu, H. Yin, and X. Wang,
\newblock Efficient, stable, and analytic differentiation of the
Sinkhorn loss,
\newblock ICLR submission, OpenReview, 2023.
\newblock OpenReview: uATOkwOZaI.

\bibitem{ThibaultEtAl2021}
A. Thibault, L. Chizat, C. Dossal, and N. Papadakis,
\newblock Overrelaxed Sinkhorn--Knopp algorithm for regularized optimal
transport,
\newblock \emph{Algorithms}, 14(5):143, 2021.
\newblock DOI: 10.3390/a14050143.

\end{thebibliography}

\end{document}
